\chapter{Two-Stage Selection with Exact Inheritance on a Wishart Family}
\label{ch:conditional-formation}

This chapter studies the reciprocal law potential of \cref{def:dyn-potential}, with its exact Gibbs rows (\cref{prop:dyn-gibbs}) and unit-Fisher dynamics, on an invariant family of matrix Wishart laws. It does not concern the global or per-agent free energies of \cref{ch:dynamics}, and its conclusions are not transferred to them. Four initially singleton agents are grouped into two parents and then into their union by a declared threshold rule applied to their evolving laws. Every parent inherits the original-time dynamics, a declared joint response interface, and the collective connection and contextual tensor of its children. Every join discards a continuous distinction between model laws.

Measured against the formation standard \cref{inh:recursive-formation}, the construction is a two-stage selection with exact inheritance, not a recursive formation. The selector reads belief shapes, and the shape law $\dot\alpha_i=-2(\alpha_i-\alpha_0)$ contains neither the coupling, nor the precisions, nor the model sector, nor any other agent. \Cref{form:kappa-independence} shows that the membership history is the count-weighted centroid dendrogram of the initial belief shapes, read on a clock whose origin is set by the threshold, that it is the same for every coupling strength, and that the renormalized null and the frozen-environment control of \cref{inh:null-flow} reproduce it, event times included, on the entire admitted class. The self costs of this construction are nonzero, so the label-preserving null is a different flow: its self weights \cref{inh:null-self-row} make the shape clocks agent dependent already at zero coupling (\cref{form:lp-null}), and integrations give a different partition sequence on an open set. The verdict of interaction necessity (F5) is therefore split between the two nulls, and the clause fails; the difference against the label-preserving null is present where no agent interacts and carries no information about the interaction. The second stage forms from every counterfactual handoff in the domain of the selector, so the counterfactual clause (F6b) fails as well, whatever the null. \Cref{form:mean-field} identifies the structural cause, which every architecture with receiver-independent targets and additive moment channels shares for selectors that read relative log-likelihood ratios. The discarded model-shape allocation is dynamically inert, and the inherited contextual geometry is prepared in the initial section, is present at zero coupling, and decays exponentially (\cref{form:geometry-decay}). What the construction establishes is exact inheritance: original-time parent dynamics with the complete channel derivative, extensive response histories under adaptive protocols, and preservation of the collective connection and contextual tensor through both joins.

The architecture is declared. Children have matrix Wishart laws with common precision in each sector, product recognition, a constant model-independent evaluator, and comparison channels depending on the population moment. Individual shapes are constant over the contextual base. Responses use fresh independent samples and return extensive sums. Membership is selected by a declared threshold rather than by a threshold derived from the potential. The preparation class is open within this invariant family. Each additive quotient is available initially, and inheritance holds for every partition at every time, so inheritance does not distinguish the selected partition; the rule selects when a quotient becomes the operative description. No law is reset at a join.

\section{Matrix laws and the declared architecture}

Let both represented spaces be the cone $\mathsf S^{++}_2$ of real symmetric
positive-definite two-by-two matrices. A random matrix $X$ or $Y$ is represented
content, distinct from a contextual coordinate. Write $\mathsf P,\mathsf Q$ for
precision matrices, to distinguish them from the principal bundle $P$. For
$\nu>1/2$ and $\mathsf P\succ0$, define
\begin{equation}
 W_2(\nu,\mathsf P)(dX)
 =\frac{\det(\mathsf P)^\nu}{\Gamma_2(\nu)}
   \det(X)^{\nu-3/2}e^{-\Tr(\mathsf P X)}dX,
 \qquad
 \Gamma_2(\nu)=\sqrt{\pi}\Gamma(\nu)\Gamma(\nu-\tfrac12).
 \label{form:wishart}
\end{equation}
Here $dX=dx_{11}dx_{12}dx_{22}$. This is the central Wishart law with real shape. Its normalizer is the multivariate gamma function, and congruence by $B$ has Jacobian $\det(B)^3$ on this cone \citep[Section 2.1 and Chapter 3]{Muirhead1982}. The range $\nu>1/2$ of real shapes for which \cref{form:wishart} is a probability density is the absolutely continuous part of Gindikin's set $\{0,\tfrac12\}\cup(\tfrac12,\infty)$, whose two isolated points give singular laws \citep{Gindikin1975}; existence questions for the noncentral case are discussed by \citet{LetacMassam2017}. No Gaussian representation or integer number of degrees of freedom is required. The proof below is self-contained.

\lemmaheading{Normalization and addition}{form:addition}
The density in \cref{form:wishart} is normalized. Its mean is
$\nu\mathsf P^{-1}$. Independent variables with laws
$W_2(\nu,\mathsf P)$ and $W_2(\mu,\mathsf P)$ have sum law
$W_2(\nu+\mu,\mathsf P)$.

\begin{proof}
At $\mathsf P=I_2$, set $x_{11}=u>0$, $x_{12}=\sqrt u z$, and
$x_{22}=v+z^2$, with $v>0$ and $z\in\R$. The determinant is $uv$, the trace is
$u+v+z^2$, and the Jacobian is $\sqrt u$. The integral is
$\Gamma(\nu)\Gamma(\nu-\tfrac12)\sqrt\pi$. Substitution
$X=\mathsf P^{-1/2}Z\mathsf P^{-1/2}$ gives the general normalizer: congruence
by a positive matrix $B$ has Jacobian $\det(B)^3$ on this cone. Differentiating
the normalizing integral with respect to $\mathsf P$ gives the stated mean.

The convolution density at $S\succ0$ contains the integral
\begin{equation}
 e^{-\Tr(\mathsf P S)}
 \int_{0<X<S}\det(X)^{\nu-3/2}
                 \det(S-X)^{\mu-3/2}dX.
\end{equation}
The substitution $X=S^{1/2}US^{1/2}$ makes it a constant times
$e^{-\Tr(\mathsf P S)}\det(S)^{\nu+\mu-3/2}$. The constant is positive and
finite by Tonelli's theorem and normalization of the convolution. Normalizing
the resulting density determines it and proves the addition formula.
\end{proof}

Initially take $N=4$ agents with
\begin{equation}
 q_i=W_2(\alpha_i,\mathsf P),\qquad
 s_i=W_2(\eta_i,\mathsf Q),\qquad
 \zeta_i=q_i\otimes s_i,
 \label{form:family}
\end{equation}
where every shape exceeds $1/2$. The model shape $\eta_i$ is unrelated to a
source-row weight. Fix $\alpha_0,\eta_0>1/2$ and $\kappa>0$. The local data are
the proper reference $r_i=W_2(\eta_0,I_2)$ and the normalized constant
evaluator $G_i(m,dX)=W_2(\alpha_0,I_2)(dX)$. For any admitted finite-moment
configuration set
\begin{equation}
 M_q=\sum_i\E_{q_i}X,\qquad
 \Lambda_q=I_2+\frac{\kappa}{N}M_q,\qquad
 \mathcal K^q_{ij}[\mathbf q](k,dX)=W_2(\alpha_0,\Lambda_q)(dX).
 \label{form:channel}
\end{equation}
The model sector uses $M_s=\sum_i\E_{s_i}Y$,
$\Lambda_s=I_2+\kappa M_s/N$, and target $W_2(\eta_0,\Lambda_s)$.
Every source, including self, is allowed with prior $\pi_{ij}=1/N$ and
positive original temperatures. Each channel ignores its particular source sample $k$ but reads the current laws
through the total moment. It therefore depends on the whole configuration, not only on the two endpoint laws, as \cref{eq:dyn-peer-costs} allows. The channel is a continuous function of $M_q$ on the finite-moment configurations, so it is jointly measurable in the configuration and the input. Independent draws across receiver children specify
the joint channel experiment used when their outputs are added.

These primitive self channels are nonidentity channels by constitutive choice. With identity self channels, \cref{def:dyn-potential} would give $D^q_{ii}=0$ while the peer costs are positive away from agreement, so the Gibbs rows would depend on each receiver's peer cost and would differ from the priors. The construction requires the nonidentity choice.

All source costs in a receiver row agree, so the exact Gibbs row equals its
prior. For the same reason, promoting the source priors to state variables of the same potential would leave them fixed. The Gibbs-reduced reciprocal potential is $\mathcal E_q+\mathcal E_s$,
where
\begin{equation}
 \mathcal E_q=\sum_i\left[
 \KL(q_i\Vert W_2(\alpha_0,I_2))
 +\KL(q_i\Vert W_2(\alpha_0,\Lambda_q))\right].
 \label{form:energy}
\end{equation}
The two sectors have separate population coupling. The evaluator is constant in its model input, and this assumption is load-bearing. By \cref{eq:dyn-local-split}, a model-dependent evaluator $G_i(m,\cdot)=W_2(\alpha_0,\Gamma(m))$ would add the term $\Tr(\E_{s_i}[\Gamma(m)]X)$ to the belief force. Because the model laws differ across agents, the precisions would then cease to be common, and \cref{form:addition} would no longer apply to the parents. The sum-pushforward of independent child evaluator outputs is a kernel from the retained model sum only because the evaluator ignores the model sample; for unequal precisions $\Gamma(m_i)$ the sum is not even a Wishart law. With the constant evaluator the belief and model sectors are uncoupled, and each sector is coupled across agents only through its own population moment.
Outside the Wishart family the displayed
functional requires finite relative entropies and moments; the differential
also requires integrable log densities and log determinants, and differentiation
under the integrals along the admitted paths.

\section{The complete smooth Fisher flow}

\propositionheading{Invariant family and original-time evolution under the reciprocal potential}{form:full-flow}
With the unit Fisher metric and all channel derivatives included,
\cref{form:family} is invariant under \cref{form:energy}. In the belief sector,
\begin{align}
 C_q&=\Lambda_q+\frac{\kappa}{N}
              (M_q-N\alpha_0\Lambda_q^{-1})
       =I_2+\frac{2\kappa}{N}M_q-\kappa\alpha_0\Lambda_q^{-1},
       \label{form:complete-force}\\
 \dot\alpha_i&=-2(\alpha_i-\alpha_0),&
 \dot{\mathsf P}&=-2\mathsf P+I_2+C_q.
 \label{form:ode}
\end{align}
The model equations exchange $\alpha,\alpha_0,\mathsf P,M_q,C_q$ with
$\eta,\eta_0,\mathsf Q,M_s,C_s$. These solutions are global on
the stated family, and every regular solution, in the sense defined in the proof, with these initial laws coincides with them. Distinct initial shapes remain distinct at every finite
time.

\begin{proof}
Put $c_0=\alpha_0-3/2$. Expanding the normalized densities gives
\begin{align}
 \mathcal E_q={}&2\sum_i\int q_i\log q_idX
 -2c_0\sum_i\E_{q_i}\log\det X+\Tr M_q
 +\Tr(\Lambda_qM_q)\notag\\
 &-N\alpha_0\log\det\Lambda_q+2N\log\Gamma_2(\alpha_0).
 \label{form:expanded-energy}
\end{align}
Since $d\Lambda_q=(\kappa/N)dM_q$, the differential of the last
moment-dependent pair is $\Tr(C_qdM_q)$. Thus the full variational force,
up to a sample-independent scalar, is
\begin{equation}
 h_i(X)=2\log q_i(X)-2c_0\log\det X+\Tr[(I_2+C_q)X].
 \label{form:density-force}
\end{equation}
The direct sender variation is zero because a mass-zero variation traversing
a constant-in-input kernel gives zero. The channel variation is not zero:
it is precisely the second term in the first expression for $C_q$.
The unrestricted density equation is
$\dot q_i=-q_i(h_i-\E_{q_i}h_i)$.
Comparing the coefficients of $\log\det X$ and $X$ in this equation with
the derivative of \cref{form:wishart} gives \cref{form:ode}; normalization
fixes the remaining scalar. This proves tangency of the full density field.

For regular Wishart data with possibly different precisions the same calculation
gives $\dot{\mathsf P}_i=-2\mathsf P_i+I_2+C_q$.
Consequently
\begin{equation}
 \mathsf P_i(t)-\mathsf P_j(t)
 =e^{-2t}(\mathsf P_i(0)-\mathsf P_j(0)).
 \label{form:precision-difference}
\end{equation}
Initially common precision remains common. On that family put
$A=\sum_i\alpha_i$ and $\bar\alpha=A/N$. Then
$M_q=A\mathsf P^{-1}$ and
$\bar\alpha(t)=\alpha_0+e^{-2t}(\bar\alpha(0)-\alpha_0)$.
The precision field commutes with $\mathsf P$, so its initial eigenvectors
remain fixed and each positive eigenvalue solves
\begin{equation}
 \dot p=f(p,\bar\alpha(t))
 =-2p+2+\frac{2\kappa\bar\alpha(t)}p
       -\frac{\kappa\alpha_0p}{p+\kappa\bar\alpha(t)}.
 \label{form:eigenvalue}
\end{equation}
Let $m_-=\min(\bar\alpha(0),\alpha_0)$ and
$m_+=\max(\bar\alpha(0),\alpha_0)$. For
$\varepsilon=\min(1/2,m_-/\alpha_0)$, $f>0$ when $0<p\le\varepsilon$.
For $L=2+\kappa m_+$, $f\le-2$ when $p\ge L$.
Every eigenvalue therefore stays between
$\min(p(0),\varepsilon)>0$ and $\max(p(0),L)<\infty$.
Shapes stay between their initial values and their reference shape, all
strictly above $1/2$. The smooth finite-dimensional field consequently has
a unique global solution on this invariant family.

Slightly enlarged compact parameter neighborhoods dominate polynomial moments
and logarithmic determinant derivatives in \cref{form:wishart}. They justify
finite KL, square-integrable scores, and the differentiations above.

Uniqueness is stated for the following class. A regular solution on $[0,\tau)$ is a family of Lebesgue densities $q_i(t,X)>0$ on $\mathsf S^{++}_2$ such that each $q_i(t)$ is a probability density with $\E|X|$, $\E|\log\det X|$, and $\E|\log q_i|$ finite, $t\mapsto M_q(t)$ is locally integrable, and for almost every $X$ the map $t\mapsto\log q_i(t,X)$ is absolutely continuous with $\partial_t\log q_i=-h_i+\E_{q_i}h_i$ for almost every $t$, where $h_i$ is \cref{form:density-force} and $t\mapsto\E_{q_i}h_i$ is locally integrable. Put $C(t)=C_q(M_q(t))$, which is locally integrable because $\Lambda_q^{-1}\preceq I_2$. For almost every $X$ the equation is scalar and linear, so
\begin{equation*}
 \log q_i(t,X)=e^{-2t}\log q_i(0,X)+(1-e^{-2t})c_0\log\det X-\Tr(\widehat C(t)X)+k_i(t),
\end{equation*}
with $\widehat C(t)=\int_0^te^{-2(t-u)}(I_2+C(u))du$ and $k_i$ a function of $t$ alone. With Wishart initial laws this is $\det(X)^{a_i(t)-3/2}e^{-\Tr(\Pi_i(t)X)}$ up to normalization, with $a_i(t)=\alpha_0+e^{-2t}(\alpha_i(0)-\alpha_0)>1/2$ and $\Pi_i(t)=e^{-2t}\mathsf P_i(0)+\widehat C(t)$. Normalizability forces $\Pi_i(t)\succ0$: if $\Pi_i(t)$ had an eigenvalue $\pi_2\le0$, then after a rotation to $\operatorname{diag}(\pi_1,\pi_2)$ the chart of \cref{form:addition} gives a factor $\int_0^\infty v^{a_i-3/2}e^{-\pi_2v}dv=\infty$. Hence $q_i(t)=W_2(a_i(t),\Pi_i(t))$, so $M_q(t)=\sum_ja_j(t)\Pi_j(t)^{-1}$ and $C$ are continuous, $\Pi_i$ is continuously differentiable, and $(a,\Pi)$ solves the closed parameter system above, whose field is smooth on $\{a>1/2,\ \Pi\succ0\}$. Uniqueness for that system identifies it with the solution constructed above. Bounded log-density charts in the sense of \cref{def:dyn-chart-flow} do not contain this family, since $\log\det X$ and $\Tr(\mathsf PX)$ are unbounded, so \cref{prop:dyn-chart-flow} is not used; the chapter works on the finite-dimensional invariant family together with this uniqueness statement.
Finally, shape differences equal their initial differences times $e^{-2t}$.
\end{proof}

The log-density equation has the form $\partial_t\log q_i=-2\log q_i-V_t+c_i(t)$, with $V_t$ common to all agents. For a fixed target this is the Fisher--Rao gradient flow of a relative entropy, whose explicit solution interpolates geometrically between the initial law and the target \citep[Appendix B]{ChenEtAl2023}; the rate is two here because two relative-entropy terms contain $q_i$. In the present construction the target moves through $M_q$, and the Wishart family carries the solution.

The precision flow is coupled across agents. For a symmetric moment
variation $H$,
\begin{equation}
 DC_q[H]=\frac{2\kappa}{N}H
 +\frac{\kappa^2\alpha_0}{N}\Lambda_q^{-1}H\Lambda_q^{-1}.
 \label{form:coupling}
\end{equation}
Its trace pairing with nonzero $H$ is strictly positive. Varying another
agent's shape produces $H=\mathsf P^{-1}d\alpha_j$ and changes the coefficient of the linear statistic $X$ in the force on $q_i$, hence the precision velocity. It does not change any shape velocity: in \cref{form:density-force} the coupling enters only the coefficient of $X$, while the coefficient of $\log\det X$ after substitution of $\log q_i$ is $2(\alpha_i-\alpha_0)$ for every $\kappa\ge0$. The coupling comes from the declared population moment;
the source rows remain at their priors and do not select preferred neighbors.

\section{A common relaxation clock}

The independence of the shape law from the interaction is an instance of a general identity for receiver-independent targets.

\propositionheading{Shared targets impose a common relaxation clock}{form:mean-field}
Consider the reciprocal potential for agents with common state and model spaces. Suppose every agent has the same model reference $r$, the same evaluator $G(m,dk)=g(dk)$, constant in $m$, and that every belief channel, for every receiver, source, and input, equals one normalized law $\nu^q[M_q]$ depending smoothly on $M_q=\sum_i\E_{q_i}T$ for a fixed measurable statistic $T:\mathsf K\to\R^d$; the model sector is analogous. Let all source priors and temperatures be positive. Along any solution of the unit-Fisher density equation of the Gibbs-reduced potential in the sense of the regular class of the proof of \cref{form:full-flow}, with $T$ in place of $X$ and with differentiation under the integrals assumed, every pair of agents satisfies
\begin{equation}
 \partial_t\log\frac{q_i}{q_j}=-2\log\frac{q_i}{q_j}+c_{ij}(t),
 \label{form:common-clock}
\end{equation}
with $c_{ij}$ independent of the sample. Hence $\log(q_i/q_j)(t)=e^{-2t}\log(q_i/q_j)(0)$ up to a sample-independent function of $t$, for every target map $\nu^q$, in particular for every coupling strength and for a constant target. If the laws lie in an exponential family with natural parameter $\theta$ and affinely independent sufficient statistic, then $\theta_i(t)-\theta_j(t)=e^{-2t}(\theta_i(0)-\theta_j(0))$.

\begin{proof}
All costs in row $i$ equal $D_i=\KL(q_i\Vert\nu^q[M_q])$, so by \cref{prop:dyn-gibbs} the Gibbs row equals the prior and $-T_q\log Z^q_i=D_i$. By \cref{eq:dyn-local-split} with a constant evaluator, the part of the reduced potential that depends on the belief laws is $\sum_l[\KL(q_l\Vert g)+\KL(q_l\Vert\nu^q[M_q])]$. By \cref{eq:dyn-kl-differential}, a variation of $q_i$ contributes $\log(q_i/g)+\log(q_i/\nu^q[M_q])$ through the two terms in which $q_i$ is the receiver, and $\langle\Psi_t,T\rangle$ through $M_q$, where $\Psi_t=-\sum_l\E_{q_l}\partial_M\log\nu^q[M]$ at $M=M_q$. Since $M_q$ is additive and every agent enters it through the same statistic, $\Psi_t$ is the same for every $i$. Thus $h_i=2\log q_i+V_t$ up to a constant, with $V_t=-\log g-\log\nu^q[M_q]+\langle\Psi_t,T\rangle$ common to all agents. The density equation gives $\partial_t\log q_i=-2\log q_i-V_t+\E_{q_i}h_i$, and subtraction gives \cref{form:common-clock}. Integrating the linear equation proves the second statement. In an exponential family, $\log(q_i/q_j)=\langle\theta_i-\theta_j,S\rangle-\psi(\theta_i)+\psi(\theta_j)$, and affine independence of $S$ identifies the coefficients.
\end{proof}

For the Wishart family the natural parameter is $(\alpha-3/2,-\mathsf P)$, so \cref{form:mean-field} contains both the shape law and \cref{form:precision-difference}. On finite positive tables the same identity is case (b) of \cref{inh:common-clock}, which also records what happens along the nulls of \cref{inh:null-flow}: with self in every source set the self costs are nonzero, the two nulls differ, and neither keeps a common clock in general. In the present family only the shape coordinates of the renormalized null and of the frozen-environment control keep the clock, because the frozen moment enters only the coefficient of $X$ (\cref{form:kappa-independence}). Any relabeling-equivariant selector that reads only centered log-likelihood ratios between agents, for example natural-parameter differences, sees along the full flow a rescaled copy of the initial data on a clock shared by every interaction in this class, so its partition sequence does not depend on the coupling. The conclusion is relative to that selector class. The cost of agent $i$ satisfies, at common precision,
\begin{equation}
 D_i-D_j=F(\alpha_i)-F(\alpha_j)+(\alpha_i-\alpha_j)\tau(t),\qquad D_i=\KL(q_i\Vert W_2(\alpha_0,\Lambda_q)),
 \label{form:cost-difference}
\end{equation}
with $F(a)=(a-\alpha_0)\psi_2(a)-\log\Gamma_2(a)$, $\psi_2=(\log\Gamma_2)'$, and $\tau=\Tr(\Lambda_q\mathsf P^{-1})-2$, which depends on $\kappa$ through $\mathsf P$ and $\Lambda_q$. A threshold rule on $|D_i-D_j|$ can therefore be coupling sensitive although every target is receiver independent. This is one reason the formation standard fixes the selector family before the flows are analyzed. Whether exact sum closure on an open fiber of child allocations is compatible with receiver-selective rows is not settled here.

A second obstruction concerns exact additive closure itself.

\propositionheading{Scalar sum closure forces a common clock}{form:scalar-closure}
Let scalar parameters $\theta_i$ in an open interval $\Theta$ evolve by $\dot\theta_i=G(\theta_i,g(t))$, where $g(t)$ is a state common to all agents and $G$ is continuous in its first argument. Suppose that for every value of $g$ the sum $G(x,g)+G(y,g)$ depends on $(x,y)\in\Theta^2$ only through $x+y$, as exact intertwining of sum retention for every two-element block requires. Then $G(\theta,g)=a(g)\theta+c(g)$, and every difference of count-weighted centroids of disjoint blocks is multiplied by the common factor $\exp\int_0^ta(g(u))du$. The ordering of all centroid distances is therefore that of the initial data at every time.

\begin{proof}
Fix $g$ and write $G(\theta)$. For $x,y\in\Theta$ the midpoint $m=(x+y)/2$ lies in $\Theta$ and $x+y=m+m$, so $G(x)+G(y)=2G(m)$. Choose $u<v$ in $\Theta$, and let $\ell$ be the affine function agreeing with $G$ at $u$ and $v$. Then $G-\ell$ satisfies the same midpoint equation and vanishes at $u$ and $v$, hence at every dyadic point of $[u,v]$, hence on $[u,v]$ by continuity. Since $u,v$ are arbitrary, $G$ is affine on $\Theta$. A count-weighted centroid $\theta_I=n_I^{-1}\sum_{i\in I}\theta_i$ then obeys $\dot\theta_I=a\theta_I+c$, so centroid differences obey a linear homogeneous equation with the common coefficient $a(g(t))$.
\end{proof}

The Wishart shapes are the instance $a=-2$, $c=2\alpha_0$, in which even $a$ is independent of the interaction. The same Jensen argument, with its consequence for threshold selectors, is \cref{inh:scalar-additive}. Vector parameters with a population-dependent generator that is not a multiple of the identity escape \cref{form:scalar-closure}, because a common anisotropic propagator can change which of two differences is smaller. Within the hypotheses of \cref{form:mean-field} this escape is closed: natural-parameter differences obey a common propagator equal to $-2$ times the identity, whatever the dimension and the coupling. An anisotropic propagator therefore needs receiver- or source-dependent targets, unequal mobilities, or a selector that reads more than natural-parameter differences; a smooth reciprocal-potential instance with exact closure is not constructed here.

\section{Parent laws and exact joining}

\theoremheading{Additive parents with the original Fisher dynamics}{form:parent}
For any disjoint current partition, retain
\begin{equation}
 n_I=|I|,\qquad A_I=\sum_{i\in I}\alpha_i,\qquad
 B_I=\sum_{i\in I}\eta_i,\qquad
 q_I=W_2(A_I,\mathsf P),\quad s_I=W_2(B_I,\mathsf Q).
 \label{form:parent-state}
\end{equation}
The parent evaluator, reference, and comparison channels are the corresponding
sum-pushforwards of independent child outputs. Their shapes are
$n_I\alpha_0,n_I\eta_0$, and their source priors are
$\pi^c_{IJ}=n_J/N$. With the primitive $N=\sum_I n_I$ kept fixed, their
ordinary unit-Fisher dynamics is
\begin{equation}
 \dot A_I=-2(A_I-n_I\alpha_0),\qquad
 \dot B_I=-2(B_I-n_I\eta_0),
 \label{form:parent-ode}
\end{equation}
together with the same precision equations \cref{form:ode}. Joining disjoint
blocks adds $(n,A,B)$ and keeps the common precisions. It commutes with laws,
channels, and evolution, and is associative.

\begin{proof}
Use the fixed sample map
$(X_i,Y_i)_{i\in I}\mapsto(\sum_iX_i,\sum_iY_i)$.
\Cref{form:addition} gives \cref{form:parent-state}. The same map on the
independent reference or channel outputs gives
$W_2(n_I\alpha_0,I_2)$ or $W_2(n_I\alpha_0,\Lambda_q)$, and the model exchange.
The parent evaluator is constant in its model input, just as the fine
evaluator is; its descent through the sum map uses this constancy. Source normalization gives $\pi^c_{IJ}=n_J/N$ by summing
primitive sender weights. Its row is again the exact Gibbs prior.

Total moments satisfy
$M_q=\sum_I\E_{q_I}X_I$ and $M_s=\sum_I\E_{s_I}Y_I$.
The parent reciprocal energy in the belief sector is
\begin{equation}
 \mathcal E_q^c=\sum_I\left[
 \KL(q_I\Vert W_2(n_I\alpha_0,I_2))
 +\KL(q_I\Vert W_2(n_I\alpha_0,\Lambda_q))\right].
\end{equation}
Its $\Lambda_q$ terms are still
$\Tr(\Lambda_qM_q)-N\alpha_0\log\det\Lambda_q$ because $\sum_I n_I=N$.
Its complete moment derivative is therefore the same $C_q$. The density-force
calculation in \cref{form:density-force}, with reference shape $n_I\alpha_0$,
gives \cref{form:parent-ode} and the original precision equation. This is the
projected fine vector field at the same time, with no mobility correction.

Adding two parents preserves total moments, counts, and forcing. Shape
equations add, and sample addition gives the joined law and normalized
channel. Associativity follows from addition of the three retained scalars
and of the sample matrices.
\end{proof}

The theorem holds for every partition at every time, so exact inheritance does not single out the partitions that the selector of the next section produces; selection and inheritance are logically independent here.

The parent is type closed in the sense of \cref{inh:type-closed}: its descended field is the unit-Fisher flow of a reciprocal potential of the parent laws, with normalized law-dependent channels and a constant evaluator. A shorter route to \cref{form:full-flow} and \cref{form:parent} uses natural coordinates. In $x=(\nu,\mathsf P)$ the law $W_2(x)$ has log density $\nu\log\det X-\Tr(\mathsf PX)-F(x)$ against $\det(X)^{-3/2}dX$, with $F(\nu,\mathsf P)=\log\Gamma_2(\nu)-\nu\log\det\mathsf P$. Its Fisher matrix is the Hessian of $F$, and $\KL(W_2(x)\Vert W_2(y))$ is the Bregman divergence of $F$ \citep{AmariNagaoka2000}, so the natural gradient of $x\mapsto\KL(W_2(x)\Vert W_2(y))$ is $x-y$. The natural gradient of a function $\Phi$ of the mean $m=\E X=\nu\mathsf P^{-1}$ with $d\Phi=\Tr(G dm)$ is $(0,-G)$. Summing these for the reference, for the target at fixed $\Lambda_q$, and for the moment dependence with $G=(\kappa/N)(M_q-N\alpha_0\Lambda_q^{-1})$ gives \cref{form:ode} at every state with individual precisions; the same sum for the parents, with targets of shape $n_I\alpha_0$, gives \cref{form:parent-ode}. The retention is not horizontally conformal for the Fisher metrics, so \cref{inh:gradient-inheritance} does not give this descent (\cref{inh:type-closure-open}).

The final single parent still has its inherited nonidentity internal-source
channel. Replacing it by an identity self-comparison would delete the second
KL term and change its flow. A parent source label is the pushforward of one
primitive source-label draw; a batch exposing separate primitive labels would
require their joint mark law. Those labels are not exposed by the interface
below.

Exact evolution does not assert equality of fine and parent energies.
For example, at $\mathsf P=I_2$ replace two shapes $\alpha_0,\alpha_0$ by
$\alpha_0+\epsilon,\alpha_0-\epsilon$. The total moment and parent state are
unchanged, but the fine energy increases by
$2\KL(W_2(\alpha_0+\epsilon,I_2)\Vert W_2(\alpha_0,I_2))+
2\KL(W_2(\alpha_0-\epsilon,I_2)\Vert W_2(\alpha_0,I_2))>0$.
The vector-field calculation, rather than energy factorization, proves
inheritance.

\section{Selected membership}

\definitionheading{Threshold membership rule}{form:threshold-rule}
Fix a threshold $\delta>0$ and begin with singletons whose initial belief shapes are pairwise farther than $\delta$ apart. For a current block read its mean belief-shape deviation
\begin{equation}
 b_I=A_I/n_I-\alpha_0,\qquad
 \dot b_I=-2b_I,\qquad
 b_{I\cup J}=\frac{n_Ib_I+n_Jb_J}{n_I+n_J}.
 \label{form:selector}
\end{equation}
Join two current blocks when their distance $|b_I-b_J|$ first reaches $\delta$ from
above. Simultaneous threshold edges join their connected components.
Recompute distances after each event.

This rule reads only retained laws and
counts, is invariant under relabeling, and contains no supplied membership
tree. The threshold specifies a persistence criterion; it is additional
constitutive data. The separation hypothesis belongs to the domain of the rule: two shapes initially closer than $\delta$ never reach $\delta$ from above and are never joined to each other, even when they are nearest neighbors. Because every score is explicit, the event times below are computed in closed form, and no discretization or event-location scheme is involved. For the report against the formation standard the selector family is this rule with $\delta\in\mathsf C=[9/100,11/100]$, declared with the constitutive data before the flows are analyzed, and the admitted class consists of the preparations of \cref{form:selection} whose initial gaps all exceed $11/100$.

\propositionheading{Coupling-independent centroid dendrogram and the controls that reproduce it}{form:kappa-independence}
On the invariant family of \cref{form:full-flow}, let $\kappa\ge0$, all precisions, and all model data be arbitrary.
\begin{enumerate}
\item[(i)] Every shape obeys $\alpha_i(t)=\alpha_0+e^{-2t}(\alpha_i(0)-\alpha_0)$, so every block score obeys $b_I(t)=e^{-2t}b_I(0)$, where $b_I(0)$ is the count-weighted centroid of the initial deviations $\alpha_i(0)-\alpha_0$, $i\in I$.
\item[(ii)] Let $N\ge2$ initial belief shapes be pairwise farther than $\delta$ apart, and suppose that at every stage of the agglomeration below the closest pair of current centroids is unique. Then \cref{form:threshold-rule} makes $N-1$ binary joins and coincides with count-weighted centroid agglomerative clustering of the initial shapes: at stage $k$ the two current blocks whose initial centroids are closest, at distance $D_k$, join at time $t_k=\tfrac12\log(D_k/\delta)$. The heights $D_k$ increase strictly; at each event every other current distance exceeds $\delta$; and the scores of the two blocks joined at $t_k$ differ by $\delta e^{-2(t-t_k)}<\delta$ at every later time. The membership tree depends only on the initial belief shapes. It is independent of $\kappa$, of the precisions, of the model sector, and of $\delta$, which shifts every event time by $-\tfrac12\log\delta$.
\item[(iii)] Coupling-deformation control. For the one-parameter family of constructions indexed by $\kappa\ge0$, whose member at $\kappa=0$ has channel targets equal to the reference and hence no law-dependent channel, the membership history $H(x)$, event times included, is the same for every $\kappa\ge0$ at every point $x$ of the admitted class. This control is not one of the controls of \cref{inh:null-flow}.
\item[(iv)] Renormalized null. The null $\Phi^0_{\mathrm{rn}}$ of \cref{inh:null-flow}, in which the peer costs and their labels are deleted, the self row has weight one, and each self channel is evaluated with the peer laws held at their initial values, exists globally on the Wishart family with individual precisions and has the same shape trajectories as $\Phi$. Hence $H^0_{\mathrm{rn}}(x)=H(x)$, event times included, and $P^0_{\mathrm{rn}}(x)=P(x)$ at every point $x$ of the admitted class, for every $\kappa\ge0$ and every $\delta$. This null is not the flow at $\kappa=0$: it keeps $\kappa$ inside each self channel, so its precisions differ and need not stay common.
\item[(v)] Frozen-environment control. The control $\Phi^{\mathrm{fr}}$ of \cref{inh:null-flow} exists globally on the same family and has the same shape trajectories, so $H^{\mathrm{fr}}(x)=H(x)$ and $P^{\mathrm{fr}}(x)=P(x)$ at every point of the admitted class.
\end{enumerate}
After a handoff every control remains the fine flow, and the selector reads the block shape sums $A_I$ computed from the fine shapes. Along $\Phi^0_{\mathrm{rn}}$ and $\Phi^{\mathrm{fr}}$ with $\kappa>0$ the precisions are individual, the block laws are not Wishart, and \cref{form:parent} and \cref{form:responses} are not claimed along them.

\begin{proof}
(i) Substituting $\log q_i=(\alpha_i-\tfrac32)\log\det X-\Tr(\mathsf PX)+\mathrm{const}$ into \cref{form:density-force} gives $h_i=2(\alpha_i-\alpha_0)\log\det X+\Tr[(I_2+C_q-2\mathsf P)X]+\mathrm{const}$. The coupling, the precision, and the other agents enter only the coefficient of $X$; matching the coefficient of $\log\det X$ gives the shape law for every $\kappa\ge0$, and at $\kappa=0$ the channel target equals the reference. Block scores are count-weighted averages of agent scores, and averaging commutes with the common factor $e^{-2t}$.

(ii) By (i), every current distance is $De^{-2t}$ with $D$ a distance between initial centroids, so a pair reaches $\delta$ at $\tfrac12\log(D/\delta)$, and the first pair to reach it is the pair with the smallest $D$. Since all initial distances exceed $\delta$, every pair reaches $\delta$ from above. Suppose $I$ and $J$ are the closest current pair, at distance $D$, and let $K$ be any other current block. The merged centroid lies strictly between the centroids of $I$ and $J$. The centroid of $K$ does not lie between them, since otherwise it would be closer than $D$ to one of them. If it lies beyond the centroid of $J$, its distance to the merged centroid is its distance to $J$ plus $n_ID/(n_I+n_J)$, which exceeds $D$; the other side is symmetric. Thus every distance after the join exceeds $D$, which proves that heights increase strictly, that centroids remain distinct, and that at the event time every other current distance is above $\delta$, so there are no cascades and every later pair again reaches $\delta$ from above. The two joined blocks keep their scores along the fine flow, so their difference is $\delta e^{-2(t-t_k)}$. This is the agglomeration that joins the closest pair of current centroids and replaces it by its count-weighted centroid. The heights $D_k$ are functions of the initial shapes alone.

(iii) is (i) and (ii), whose conclusions do not depend on $\kappa$.

(iv) The renormalized null deletes the peer costs with their labels and evaluates each self channel with the peer laws held at their initial values. Its belief potential is $\sum_i[\KL(q_i\Vert W_2(\alpha_0,I_2))+\KL(q_i\Vert W_2(\alpha_0,\Lambda_i^0))]$ with $\Lambda_i^0=I_2+(\kappa/N)(\E_{q_i}X+M_{-i}^0)$, where $M_{-i}^0$ is the initial moment of the other agents. The calculation of \cref{form:full-flow} with this single-agent moment gives the same shape law and $\dot{\mathsf P}_i=-2\mathsf P_i+I_2+C_i^0$, where $C_i^0=\Lambda_i^0+(\kappa/N)(\alpha_i\mathsf P_i^{-1}-\alpha_0(\Lambda_i^0)^{-1})$. For a unit eigenvector $v$ of the smallest eigenvalue $p$ of $\mathsf P_i$, $\Lambda_i^0\succeq I_2$ gives $v^T\dot{\mathsf P}_iv\ge-2p+2-\kappa\alpha_0/N+2\kappa\alpha_i/(Np)$, which is positive for small $p$ because $\alpha_i$ stays above $\min(\alpha_i(0),\alpha_0)>1/2$. With the smallest eigenvalue bounded below, the largest eigenvalue $p$ satisfies $\dot p\le-2p+K$ for a constant $K$ on each finite interval. The precisions therefore remain in compact subsets of the cone, and the solution is global and unique as in \cref{form:full-flow}. The selector reads only belief shapes, whose trajectories are the same under both flows, also after every handoff, so the membership histories agree, event times included.

(v) Along $\Phi^{\mathrm{fr}}$ nothing is deleted, every row stays at its prior because all costs in a row agree, and in every term that enters the force on $q_i$ the other laws take their initial values. The force is $h_i=2\log q_i-2c_0\log\det X+\Tr[(I_2+C_i^{\mathrm{fr}})X]$ with $C_i^{\mathrm{fr}}=\Lambda_i^0+(\kappa/N)(\alpha_i\mathsf P_i^{-1}+M_{-i}^0-N\alpha_0(\Lambda_i^0)^{-1})$, so the shape law is that of (i) and $\dot{\mathsf P}_i=-2\mathsf P_i+I_2+C_i^{\mathrm{fr}}$. At the smallest eigenvalue $p$ of $\mathsf P_i$, $v^T\dot{\mathsf P}_iv\ge-2p+2-\kappa\alpha_0+\kappa\alpha_i/(Np)$, which is positive for small $p$ when $\kappa>0$, and the case $\kappa=0$ is the flow $\dot{\mathsf P}_i=-2\mathsf P_i+2I_2$; the largest eigenvalue is bounded as in (iv). The solution is global, and the selector sees the shape trajectories of (i).
\end{proof}

The label-preserving null of \cref{inh:null-flow} keeps the peer labels and priors, so with Gibbs rows the self row of receiver $i$ has the weight \cref{inh:null-self-row}, which is below one.

\propositionheading{The label-preserving null has agent-dependent clocks}{form:lp-null}
Along $\Phi^0_{\mathrm{lp}}$, with self prior $a=1/N$ and belief temperature $T_q$, put $D_{ii}=\KL(q_i\Vert W_2(\alpha_0,\Lambda_i^0))$, with $\Lambda_i^0$ and $C_i^0$ as in the proof of \cref{form:kappa-independence}, and let $\beta_i=ae^{-D_{ii}/T_q}/(ae^{-D_{ii}/T_q}+1-a)\in(0,1/N]$. The Wishart family with individual precisions is invariant, with
\begin{equation}
 \dot\alpha_i=-(1+\beta_i)(\alpha_i-\alpha_0),\qquad
 \dot{\mathsf P}_i=-(1+\beta_i)\mathsf P_i+I_2+\beta_iC_i^0,
 \label{form:lp-null-law}
\end{equation}
and the model analog. For $\kappa\alpha_0\le N^2$ the solution exists for all $t\ge0$ and is unique on this family. At $\kappa=0$ and $\mathsf P_i(0)=I_2$ the precisions stay at $I_2$ and each deviation $b_i=\alpha_i-\alpha_0$ obeys the autonomous equation
\begin{equation}
 \dot b=-(1+\beta(b))b,\qquad
 \beta(b)=\frac{1}{1+(N-1)e^{F_0(b)/T_q}},\qquad
 F_0(b)=\KL(W_2(\alpha_0+b,I_2)\Vert W_2(\alpha_0,I_2)),
 \label{form:lp-null-deviation}
\end{equation}
so that agents farther from $\alpha_0$ relax more slowly, and the clocks differ between agents whenever their shapes differ in distance from $\alpha_0$.
\begin{proof}
With peer costs set to zero and exact Gibbs rows, the row term of receiver $i$ is $-T_q\log(ae^{-D_{ii}/T_q}+1-a)$, whose first variation in $q_i$ is $\beta_i$ times that of $D_{ii}$. The first variation of $D_{ii}$ is $\log q_i-\log W_2(\alpha_0,\Lambda_i^0)$ plus the channel term $\Tr[(\kappa/N)(\alpha_i\mathsf P_i^{-1}-\alpha_0(\Lambda_i^0)^{-1})X]$, so the force is $(1+\beta_i)\log q_i-(1+\beta_i)c_0\log\det X+\Tr[(I_2+\beta_iC_i^0)X]$ up to a constant, and matching coefficients as in \cref{form:full-flow} gives \cref{form:lp-null-law}. The closed form $\KL(W_2(a,\mathsf P)\Vert W_2(b,\mathsf Q))=(a-b)\psi_2(a)-\log\Gamma_2(a)+\log\Gamma_2(b)+b\log\det(\mathsf P\mathsf Q^{-1})+a(\Tr(\mathsf Q\mathsf P^{-1})-2)$ makes $\beta_i$ a smooth function of the state. Shapes move monotonically from $\alpha_i(0)$ toward $\alpha_0$. At the smallest eigenvalue $p$ of $\mathsf P_i$ with unit eigenvector $v$, $\Lambda_i^0\succeq I_2$ gives $v^TC_i^0v\ge1-\kappa\alpha_0/N$, so, since $0<\beta_i\le1/N$, $v^T\dot{\mathsf P}_iv\ge-2p+1+\beta_i(1-\kappa\alpha_0/N)\ge-2p+1+1/N-\kappa\alpha_0/N^2$ when $\kappa\alpha_0\geq N$ and $\ge-2p+1$ otherwise. For $\kappa\alpha_0\le N^2$ both bounds are positive for small $p$. With the smallest eigenvalue bounded below, $C_i^0$ is bounded and the largest eigenvalue satisfies $\dot p\le-p+K$. The parameter field is smooth on the invariant region, which gives global existence and uniqueness. At $\kappa=0$, $\Lambda_i^0=C_i^0=I_2$ and $\dot{\mathsf P}_i=(1+\beta_i)(I_2-\mathsf P_i)$, so $\mathsf P_i\equiv I_2$ from $I_2$, and then $D_{ii}=F_0(b_i)$, which is zero at $b=0$ and increases with $|b|$ because $F_0$ is a relative entropy in the shape of an exponential family with convex log partition. Hence $\beta$ decreases in $|b|$.
\end{proof}

As $T_q\to\infty$ every $\beta_i$ tends to $1/N$, the clocks become the common rate $1+1/N$, and the partition sequence of the full flow returns. At $T_q=1$ the agent-dependent clocks change the partition sequence. For $\alpha_0=1$, $\kappa=0$, $\mathsf P(0)=I_2$, and initial shapes $(1.2807,1.5057,1.8665,2.3057)$, which satisfy \cref{form:gaps} for every $\delta\in\mathsf C$, the full flow at $\delta=1/10$ forms $\{1,2\}$, $\{3,4\}$, and their union at $t=0.406$, $0.740$, $0.968$, while integration of \cref{form:lp-null-law} with the selector applied to fine shape sums forms $\{1,2\}$, $\{1,2,3\}$, and the union at $t=0.689$, $1.340$, $1.776$. These sequences differ even modulo exchange of disjoint merges. At each null event every other block distance exceeds $\delta$ by at least $3.5\times10^{-3}$, and every crossing is transverse, at five values of $\delta$ spanning $\mathsf C$ with $\kappa=0$; the null sequence is the same at $\kappa=0.05$ and $0.1$, and the same two sequences were obtained under twelve random perturbations of size $4\times10^{-3}$ in the shapes and $10^{-2}$ in $\mathsf P(0)$, with $\kappa\in\{0,0.01\}$ and $T_q\in[0.9,1.1]$. By continuous dependence of both flows on the initial data, $P^0_{\mathrm{lp}}\neq P$ then holds on a nonempty open subset of the admitted class, already at zero coupling. This conclusion rests on floating-point integration with these margins, not on a validated enclosure. At $\kappa=T_q=1$, $\mathsf P(0)=I_2$, and the example shapes below, the label-preserving null forms the same blocks as the full flow, at the later times $0.650$, $1.111$, $2.058$. Where $\kappa\alpha_0>N^2$ the existence of this null is not shown, and (F5) against it is undefined there.

Centroid clustering is one of the pair-group strategies of numerical taxonomy \citep{SokalMichener1958}, unified with the other agglomerative strategies by \citet{LanceWilliams1967}. In one dimension it has no reversals, which is the monotonicity in (ii). The rule is not single-linkage clustering. For initial deviations $0,1,2.05,3.5$, single linkage joins the third point to the first pair before the last point, whereas the rule forms the two pairs first, at heights $1$ and $1.45$, and their union at $2.275$. The union time in \cref{form:second-time} below already carries the centroid distance $d_2+(d_1+d_3)/2$ rather than the single-linkage distance $d_2$. The construction also differs from bounded-confidence dynamics \citep{HegselmannKrause2002}, in which the interaction determines which agents cluster. What it adds to centroid clustering is exact inheritance of laws, dynamics, responses, and geometry at each join.

\theoremheading{Two selected stages on the four-agent family}{form:selection}
Let $u_1<u_2<u_3<u_4$ be the sorted initial belief shapes, with consecutive
gaps $d_1,d_2,d_3$, and suppose
\begin{equation}
 \delta<\min(d_1,d_3)<d_2,\qquad \max(d_1,d_3)<d_2+\tfrac12\min(d_1,d_3).
 \label{form:gaps}
\end{equation}
All four shapes exceed $1/2$, while model shapes vary independently above
$1/2$. Under \cref{form:threshold-rule}, the two lower and the two upper shapes
are joined into parents at
\begin{equation}
 t_L=\tfrac12\log(d_1/\delta),\qquad
 t_R=\tfrac12\log(d_3/\delta).
 \label{form:first-times}
\end{equation}
Their union is joined at
\begin{equation}
 t_2=\tfrac12\log\left(\frac{d_2+(d_1+d_3)/2}{\delta}\right)
 >\max(t_L,t_R).
 \label{form:second-time}
\end{equation}
There are no intervening cross joins or cascades. Both first parents coexist
on $[\max(t_L,t_R),t_2)$. Each joined child-block distance remains below
threshold afterward, and all fine belief laws remain distinct at every finite
time. When strict inequalities replace ties, \cref{form:gaps} is also necessary for this membership history. Each join forgets a continuous model-shape distinction even if all
membership labels and event times are retained.

\begin{proof}
Apply \cref{form:kappa-independence}. Suppose $d_1\le d_3$. Since $d_1<d_2$, the lower pair is the closest initial pair. After it joins, the distance of its centroid to $u_3$ is $d_2+d_1/2$, which exceeds $d_3$ by \cref{form:gaps}, so the upper pair joins next. The distance between the two parent centroids is $d_2+(d_1+d_3)/2$, giving \cref{form:second-time}. The opposite order follows by exchanging upper and lower pairs. If $d_1=d_3$, exactly the two disjoint outer edges reach threshold together, and the component rule forms the same pairs. Conversely, if $d_2<\min(d_1,d_3)$ the middle pair joins first, and if $\max(d_1,d_3)>d_2+\tfrac12\min(d_1,d_3)$ the first parent absorbs the nearer middle shape before the other pair joins.

The persistence statement is part of \cref{form:kappa-independence}; it follows from the equations rather than from storing a membership flag. Primitive shape
differences never vanish at finite time. The two first-parent laws also
remain distinct, because their mean belief shapes remain different.

For loss at a first join, fix all belief shapes and perturb two initial model
shapes in that pair by $(\eta_i,\eta_j)\mapsto
(\eta_i+\epsilon,\eta_j-\epsilon)$ for sufficiently small $\epsilon$.
The parent model sum, both population moments, all subsequent precision
evolution, and every membership event and time are unchanged. Individual
model laws change continuously. Distinct fixed belief shapes distinguish
the labels, so this change is not a permutation. At the second join, transfer
model shape between the two first-parent sums while keeping their union sum
fixed. This gives a further continuous lost distinction and still leaves
the belief-only selection schedule unchanged.
\end{proof}

The discarded allocation is a dynamical spectator in both descriptions. Individual model shapes obey $\dot\eta_i=-2(\eta_i-\eta_0)$, the model precision depends on the model shapes only through their total, the belief sector never reads the model laws because the evaluator is constant, and the selector reads only belief shapes. The witness therefore establishes noninjectivity of the retention, not loss of a coordinate to which any agent, channel, or selector responds. The belief allocation that the selector reads is not lost given the event record: at a join the two child scores differ by exactly $\delta$, so the parent sum, the child counts, the join time, $\delta$, and $\alpha_0$ determine the child belief shapes up to exchange. The loss witness accordingly uses the model-shape allocation. Already observed sample records remain with the controller; they
are not erased or asserted to have the same law across different initial
preparations. Loss concerns the current parent state and the future interface
specified next.

The hierarchy is transient. All shapes converge to $\alpha_0,\eta_0$, and every precision eigenvalue converges to the unique positive zero of $f(\cdot,\alpha_0)$ or of its model analog, so every agent tends to one common law. The children are distinct at every finite time, and the consensus is asymptotic rather than reached in finite time.

The preparation class is nonempty and relatively open. For example take
$\alpha_0=\eta_0=\kappa=1$, $\delta=1/10$,
$(u_1,u_2,u_3,u_4)=(1,6/5,9/5,21/10)$, and model shapes in a neighborhood
of $(4/5,11/10,7/5,17/10)$. The gaps are $(1/5,3/5,3/10)$, and the events
occur at $\tfrac12\log2$, $\tfrac12\log3$, and $\tfrac12\log(17/2)$. The two first parents coexist for a time $\tfrac12\log(17/6)$.
A sufficiently small neighborhood and every relabeling obey
\cref{form:gaps}. Openness is within the common-precision family, not among
all positive density tuples. The smooth statistical retention applies on
each fixed membership branch; no global smoothness across switching
boundaries is required.

Against \cref{inh:recursive-formation} the report is as follows. The objective is the reciprocal potential, and the dynamics is the unit-Fisher flow of \cref{form:full-flow} on the parameter manifold of the common-precision Wishart family; no chart of the full density flow is declared, since bounded log-density charts do not contain the family. (F1) holds in the family grade: the admitted class is relatively open in the invariant family defined by Wishart laws with common precision in each sector, which has dimension $14$, infinite codimension among positive densities, and is not contained in finitely many orbits. (F2) holds for the selector family $\delta\in\mathsf C=[9/100,11/100]$, which reads only retained shapes and counts and so satisfies the readout convention, with an invariant block relation as persistence criterion. The parent protocol rule assigns to a block $I$ the fresh block-sum queries of the response section below, read through whole current blocks; the probe rule adds, before a handoff of $I$, the protocols of the blocks merged into it, which before a first join are fresh samples of single children and before the second join fresh sums over each first parent. (F3) holds for these classes by \cref{form:parent}, \cref{form:responses}, and \cref{form:geometry}. (F4) holds: the model-shape transfer in the proof of \cref{form:selection} gives two reachable post-handoff configurations with the same retained state, separated by the fresh model sample of one child, and at the second join by the fresh model sum of one first parent. The separation is transient, since every model shape relaxes to $\eta_0$ at rate two, and it is not of symmetry type, since that relaxation contracts every fiber displacement by $e^{-2t}$; the separated coordinate is moreover dynamically inert. (F5) fails against the renormalized null at every point and for every $\delta\in\mathsf C$ (\cref{form:kappa-independence} (iv)), while against the label-preserving null the partition sequence differs on an open set already at zero coupling (\cref{form:lp-null} and the integrations after it). The verdict is split, and by \cref{inh:recursive-formation} the clause fails. The selector-free discrepancy is the shape contraction rate, which is two along $\Phi$, $\Phi^0_{\mathrm{rn}}$, and $\Phi^{\mathrm{fr}}$ and the agent-dependent $1+\beta_i$ along $\Phi^0_{\mathrm{lp}}$ at every coupling, including zero; $P^{\mathrm{fr}}=P$, and the coupling-deformation control of \cref{form:kappa-independence} (iii) shows no dependence on $\kappa$ at all. (F6a) holds: the parents select their union under the inherited dynamics with the same selector, constants, and data, without a reset. For (F6b) declare $\mathcal U$ to be the whole common-precision family with shapes above $1/2$, and let $\mathcal C$ be built by \cref{inh:counterfactual-class} at the time the second first parent forms. Every state of $\mathcal C$ lies in the domain of the selector, so its two parent scores differ by more than $\delta$, and that difference contracts by $e^{-2t}$ whatever the fine allocation; the union forms from every state of $\mathcal C$, the stages form a supplied-scale sequence, and (F6b) fails, against either null. The strong grade fails with (F5); the amplifying grade is not evaluated. The construction is a two-stage selection with exact inheritance along a prepared contraction ordering, with (F5) and (F6b) failing. Since (F6b) fails independently, this classification does not depend on counting a split (F5) verdict as a failure; the (F5) verdict itself does.

\section{Exact joint response histories}

The response interface is noninvasive. At each query, draw fresh samples
independently across children and sectors from $\bigotimes_i(q_i\otimes s_i)$.
For the current disjoint partition $\Pi$, form the entire vector
\begin{equation}
 Z_\Pi=((X_I,Y_I))_{I\in\Pi},\qquad
 X_I=\sum_{i\in I}X_i,\quad Y_I=\sum_{i\in I}Y_i.
\end{equation}
An action may choose any normalized readout kernel on this vector. It can
report overlapping unions of whole current blocks using the shared batch,
so no independence between such overlapping reports is assumed. Query
times and readouts may depend measurably on carried record history and
retained state. Recording an output does not reset a belief or model law.

\corollaryheading{Adaptive extensive-response inheritance}{form:responses}
For every admitted preparation and every finite sequence of finite query
times selected by the preceding protocol, the fine and parent descriptions
have the same joint law of records, actions, query times, and retained
states. This includes the membership events of \cref{form:selection}.

\begin{proof}
Independence across disjoint blocks and \cref{form:addition} give the
conditional batch kernel
\begin{equation}
 K_\Pi=\bigotimes_{I\in\Pi}
       [W_2(A_I,\mathsf P)\otimes W_2(B_I,\mathsf Q)].
 \label{form:batch}
\end{equation}
Initialize the retained process from the same fine configuration.
Condition on the retained state and carried history at a query. The same
policy chooses the next waiting time and readout in both descriptions.
\Cref{form:parent} gives equal retained states after waiting, and
\cref{form:threshold-rule} gives the same intervening events and addition handoffs.
Freshness makes the next conditional batch law exactly \cref{form:batch};
composition with the same normalized readout gives the same next-record
kernel. Induction proves equality of every finite joint history.
Independent external randomization can be conditioned on in the same proof.
\end{proof}

The result is a corollary of two facts: the interface exposes only block sums, whose law depends on the fine allocation only through $A_I$ and $B_I$, and the retained state is a deterministic function of the preparation and time. Records never feed back into laws, so adaptivity selects only query times and readouts, and the induction carries no further content. Every decision functional of the joint record law, such as the expected cost of an adaptive policy or the entropy of the record history, takes the same value in both descriptions; information gain about the discarded allocation does not transfer, because the parent has no such coordinate. The informative extension, in which records drive updates of the laws, is outside the declared interface.

The corollary compares two descriptions of each given preparation. A final
aggregate is not claimed to reconstruct previously discarded records or the
distribution of earlier records from a different preparation. Queries cutting
through a current block need additional intersection data. For example,
same-batch sums over primitive subsets $C,D$ have
\begin{equation}
 \operatorname{Cov}(\Tr(UX_C),\Tr(VX_D))
 =A_{C\cap D}\Tr(\mathsf P^{-1}U\mathsf P^{-1}V),
\end{equation}
which can expose a discarded shape sum. Repeated hidden-child sampling,
child-specific Bayesian updates, and interventions changing the reference
potential are outside the declared interface.

\section{Inherited collective geometry}

Use the trivial principal bundle $\mathcal C\times SO(2)$ and the common
left sample action $X\mapsto gXg^T$, $Y\mapsto gYg^T$. It sends precisions
to $g\mathsf P g^T,g\mathsf Qg^T$. The identity-precision references are
invariant under this action. One rotation acting on belief and model samples simultaneously is a declared choice of collective group. Because the evaluator is constant, the constitutive data are also invariant under independent rotations of the two sectors, and nothing in the potential ties them. Use the collective product-law Fisher metric
with unit weights, the common comparison frame, and the convention
$D^A p=dp+J_pA$ of \cref{ch:geometry}. Let
$\mathcal C=(a_-,a_+)\times(-\chi_0,\chi_0)$, with $0<a_-<a_+$, and
prepare the contextual section by
\begin{equation}
 \mathsf P(0;a,\chi)
 =R_\chi\operatorname{diag}(e^a,e^{-a})R_\chi^T,\qquad
 \mathsf Q(0)=\operatorname{diag}(e^{\tilde a},e^{-\tilde a}),\qquad \tilde a>0.
 \label{form:context}
\end{equation}
Here $R_\chi$ is plane rotation by $\chi$. All individual shape parameters
are constant over $\mathcal C$. Consequently the selection schedule is
also constant over this contextual chart. The flow acts pointwise in the context, and no vector field, row, channel, or record depends on the connection: the connection below is a diagnostic of the evolved sections, and its curvature is the image under the flow of the prepared contextual dependence in \cref{form:context}.

\theoremheading{Exact connection and contextual-tensor inheritance for a prepared contextual family}{form:geometry}
For every block, before and after either join, sum retention preserves its
collective minimizing connection class and minimized contextual tensor for the common $SO(2)$.
For \cref{form:context}, that tensor has rank two and the connection has
nonzero curvature at every finite time, for every $\kappa\ge0$.

\begin{proof}
Differentiating \cref{form:wishart} at fixed shape gives the precision score
and Fisher metric
\begin{equation}
 \ell_U(X)=\alpha\Tr(\mathsf P^{-1}U)-\Tr(UX),\qquad
 g^F(U,V)=\alpha\Tr(\mathsf P^{-1}U\mathsf P^{-1}V).
 \label{form:fisher}
\end{equation}
The second identity follows by differentiating the normalizing integral
twice. For common precision variations $(U,V)$, the summed fine score is
\begin{equation}
 A_I\Tr(\mathsf P^{-1}U)-\Tr(UX_I)
 +B_I\Tr(\mathsf Q^{-1}V)-\Tr(VY_I).
\end{equation}
It equals the parent score and is measurable from the retained sums.
Both contextual and common rotation-orbit directions have this form.
Their Fisher pairings and all mixed pairings are therefore preserved.
The infinitesimal stabilizers agree as well: in each description they are
the rotations commuting with both precisions. Thus the normal equations
for the minimizing connection have the same solutions modulo the same
isotropy, and the residual tensors agree. This also verifies the mixed-loss
criterion in \cref{geo:retained-connection}. The sample addition channel is
fixed, so it contributes no moving-channel derivative.

For explicit curvature, let $p_\pm(t,a)$ be the two eigenvalues of
$\mathsf P(t;a,\chi)$, and let $\tilde p_\pm(t)$ be those of $\mathsf Q(t)$.
The scalar field in \cref{form:eigenvalue} has
\begin{equation}
 f_p=-2-\frac{2\kappa\bar\alpha}{p^2}
       -\frac{\kappa^2\alpha_0\bar\alpha}{(p+\kappa\bar\alpha)^2}<0.
\end{equation}
Its solution map $\Phi_t(p_0)$ satisfies
$\partial_{p_0}\Phi_t=\exp(\int_0^t f_pdu)>0$ for finite $t$.
Hence $p_+=\Phi_t(e^a)>p_-=\Phi_t(e^{-a})>0$,
$\partial_a p_+>0$, and $\partial_a p_-<0$.
The model eigenvalues also remain distinct and have no $a$ dependence.
Define the rotational inertias
\begin{equation}
 \mathcal I_1=2A_I\frac{(p_+-p_-)^2}{p_+p_-},\qquad
 \mathcal I_2=2B_I\frac{(\tilde p_+-\tilde p_-)^2}{\tilde p_+\tilde p_-}>0.
 \label{form:inertias}
\end{equation}
Writing $r=p_+/p_->1$ gives $\mathcal I_1=2A_I(r+r^{-1}-2)$ and
$\partial_a\mathcal I_1>0$. Radial precision variations are diagonal in an
eigenframe, while orbit variations are off-diagonal, so their Fisher
pairing vanishes. The collective normal equation then gives
\begin{align}
 A^F&=-\frac{\mathcal I_1}{\mathcal I_1+\mathcal I_2}d\chi,&
 h^{\min}&=g_{aa}da^2+\frac{\mathcal I_1\mathcal I_2}{\mathcal I_1+\mathcal I_2}d\chi^2,
 \label{form:connection}\\
 g_{aa}&=A_I\left[
  \left(\frac{\partial_a p_+}{p_+}\right)^2+
  \left(\frac{\partial_a p_-}{p_-}\right)^2\right]>0.
 \notag
\end{align}
Identifying the one-dimensional Lie algebra with its rotation generator,
the curvature is
\begin{equation}
 F^F=dA^F
 =-\frac{(\partial_a\mathcal I_1)\mathcal I_2}{(\mathcal I_1+\mathcal I_2)^2}da\wedge d\chi\ne0.
 \label{form:curvature}
\end{equation}
These computations hold for every retained block and its fine constituents
with the same totals. None of them uses $\kappa>0$. Under a passive local rotation by $\vartheta(c)$,
the angular momentum one-form gains $(\mathcal I_1+\mathcal I_2)d\vartheta$, so
$A^{F\prime}=A^F-d\vartheta$. The tensor is invariant, with exactly the
affine connection law of the general construction.
\end{proof}

The minimizing connection is the mechanical connection of a two-component rotor with locked inertia $\mathcal I_1+\mathcal I_2$ \citep{Smale1970}, and $h^{\min}$ is the pullback of the quotient metric of the corresponding Riemannian submersion \citep{ONeill1966}. Its curvature measures how the ratio of belief to model rotational inertia varies along the context. Two features delimit it. First, the curvature belongs to the tied belief and model connection. For the independent group $SO(2)_q\times SO(2)_s$, which the constitutive data also admit, the Fisher cost of a connection $(\xi_q,\xi_s)$ is $\mathcal I_1(d\chi+\xi_q)^2+\mathcal I_2\xi_s^2+g_{aa}da^2$, minimized by $\xi_q=-d\chi$ and $\xi_s=0$; that connection is flat and the residual tensor $g_{aa}da^2$ has rank one. Second, the geometry is prepared rather than produced: it is present at $t=0$ and at $\kappa=0$, where the agents do not interact, and the next proposition shows that the flow erases it. At finite times $p_+>p_-$, so the infinitesimal isotropy is zero and the constant-rank hypothesis of \cref{ch:geometry} holds; agreement of the children's laws does not affect it, because the collective action is common.

\propositionheading{Decay of the inherited contextual geometry}{form:geometry-decay}
Let $p_*$ be the unique positive zero of $f(\cdot,\alpha_0)$ in \cref{form:eigenvalue}, and put $\lambda=-f_p(p_*,\alpha_0)$, so that $\lambda>2$ for $\kappa>0$ and $\lambda=2$ for $\kappa=0$. For every context, $\mathsf P(t)\to p_*I_2$, and $e^{\lambda t}(p_+-p_-)$, $e^{\lambda t}\partial_ap_+$, and $e^{\lambda t}\partial_ap_-$ converge to nonzero limits. Consequently $e^{2\lambda t}g_{aa}$ and $e^{2\lambda t}\mathcal I_1$ converge to positive limits, the minimized tensor tends to zero, and the limit laws are isotropic, with stabilizer all of $SO(2)$. For $\alpha_0=\kappa=1$, $p_*=\sqrt2$ and $2\lambda=12-4\sqrt2$.

\begin{proof}
Since $f_p<0$ and $f(p,\alpha_0)$ tends to $+\infty$ as $p\to0$ and to $-\infty$ as $p\to\infty$, the zero $p_*$ is unique. On the compact invariant region of the proof of \cref{form:full-flow}, $f$ and $f_p$ are Lipschitz. With $y=p-p_*$, write $\dot y=[f(p,\bar\alpha)-f(p_*,\bar\alpha)]+[f(p_*,\bar\alpha)-f(p_*,\alpha_0)]$. Since $f_p\le-2$, the first bracket has the sign of $-y$ and modulus at least $2|y|$, and the second is bounded by a constant times $|\bar\alpha-\alpha_0|=O(e^{-2t})$. Hence $|y(t)|\le(|y(0)|+Kt)e^{-2t}$, uniformly for initial values in a compact set. Then $f_p(p(u),\bar\alpha(u))+\lambda$ is integrable on $[0,\infty)$, so $\partial_{p_0}\Phi_t=e^{-\lambda t}(c(p_0)+o(1))$ with $c(p_0)>0$ continuous, uniformly on compact sets of $p_0$. The identities $\partial_ap_\pm=\pm e^{\pm a}\partial_{p_0}\Phi_t(e^{\pm a})$ and $p_+-p_-=\int_{e^{-a}}^{e^a}\partial_{p_0}\Phi_t(p_0)dp_0$ give the stated limits, and the formulas for $g_{aa}$ and $\mathcal I_1$, with $A_I\to n_I\alpha_0$ and $p_\pm\to p_*$, give the rest. For $\alpha_0=\kappa=1$, substitution verifies $f(\sqrt2,1)=0$ and $f_p(\sqrt2,1)=-6+2\sqrt2$.
\end{proof}

For the example constants and the block of the first two agents at $a=3/10$, $g_{aa}$ equals $4.40$, $2.69\times10^{-3}$, and $5.20\times10^{-6}$ at $t=0,1,2$; by the root join at $t_2\approx1.07$ it is below $10^{-3}$ of its initial value. The model inertia decays at the rate given by the model analog of $\lambda$. Because $A^F$ is a ratio of two vanishing inertias, the limit of the finite-time curvature depends on whether the two rates agree: for $\alpha_0=\eta_0=\kappa=1$, model log-eigenvalue $\tilde a=3/10$, and the example model shapes it tends to approximately $-1.81$, a value that depends on $\tilde a$ (about $-1.55$ at $\tilde a=0.2$ and $-1.27$ at $\tilde a=0.5$), while for $\alpha_0=1,\eta_0=2$ the connection tends to the flat value $-d\chi$ and the curvature to zero. In neither case is the limit a property of the limit laws, at which the connection is undetermined.

\remarkheading{Persistent contextual geometry from anisotropic references}{form:anisotropic-reference}
Replace the belief evaluator by $W_2(\alpha_0,\Gamma(c))$ with $\Gamma(a,\chi)=R_\chi\operatorname{diag}(e^a,e^{-a})R_\chi^T$, the belief channel target by $W_2(\alpha_0,\Lambda_q)$ with $\Lambda_q=\Gamma+\kappa M_q/N$, and the model reference by $W_2(\eta_0,\operatorname{diag}(e^{\tilde a},e^{-\tilde a}))$, with the analogous model channel. The calculation of \cref{form:full-flow} gives $C_q=\Lambda_q+(\kappa/N)(M_q-N\alpha_0\Lambda_q^{-1})$ and $\dot{\mathsf P}=-2\mathsf P+\Gamma+C_q$, with the same shape law. Precisions remain common, evaluators and targets still add under sums, and \cref{form:parent}, \cref{form:responses}, \cref{form:kappa-independence}, and the selection schedule are unchanged. With the section $\mathsf P(0;a,\chi)=\Gamma(a,\chi)$, the eigenvalues converge to the zeros of $-2p+2\gamma+2\kappa\alpha_0/p-\kappa\alpha_0p/(\gamma p+\kappa\alpha_0)$ with $\gamma=e^{\pm a}$, which differ and depend on $a$; the limit tensor and curvature are computed from these fixed points. For the example constants at $a=\tilde a=3/10$ the block tensor tends to $\operatorname{diag}(1.168,0.199)$ and the curvature to approximately $-1.689$, so rank two and nonzero curvature persist at infinite time. Both sectors need anisotropic local data; with an isotropic model reference the model sector isotropizes and the curvature tends to zero. The local data are then contextual sections that transform under the common action, so the statement above that the references are invariant no longer holds; the curvature remains a property of the declared contextual data, not of the interaction.

The context-constant shape assumption is needed only for the tensor. On every smooth section of the common-precision family the fine and retained Fisher costs of a belief-sector tangent vector $v=(\delta\alpha,U)$ differ by $\Delta(v,v)=\sum_{i\in I}\psi_2'(\alpha_i)\delta\alpha_i^2-\psi_2'(A_I)(\sum_{i\in I}\delta\alpha_i)^2$, plus the model analog, because the precision block of the Wishart Fisher matrix is linear in the shape and the mixed block adds over children. This vanishes on every vector without a shape component, in particular on the orbit directions, so by \cref{geo:retained-connection} the minimizing connection class is inherited for every such section, while by \cref{geo:minimal-contraction} the minimized tensors differ by $\Delta$ evaluated on the contextual derivatives of the shapes, which is positive unless every individual shape is context constant. For the tensor the assumption is substantive. If
$d\eta_1=1,d\eta_2=-1$ while precisions and all other shapes are fixed,
the parent model sum has zero derivative, whereas the fine score is the
centered difference of two log determinants with positive variance.
It is orthogonal to rotation scores because determinant is conjugation
invariant. Its fine minimized length is positive and its retained length
is zero. Such a contextual direction violates exact tensor inheritance.

\openproblemheading{Approximate inheritance off the common-precision family}{form:approximate-open}
The common-precision family has positive codimension, and exact sum laws fail off it. By \cref{form:precision-difference} and the shape law, the ratio of precision spread to shape spread is constant in time, so at a join the precision mismatch is $\delta$ times a preparation constant. Under moment matching, the first-order term of the log-Laplace transform of a sum of Wishart laws with nearby precisions cancels, so the sum law deviates from the Wishart family at second order in the precision spread. Prove that on an open set of Wishart preparations with unequal precisions the moment-matched parent inherits the law, the dynamics, and the extensive responses up to a stated relative-entropy defect at each join that decays afterward, and determine its order.
